
\usepackage{titlesec}
\usepackage{url}

\newsavebox{\tempbox}
\sbox{\tempbox}{\raisebox{-4.5pt}%
{\parbox{3cm}{\textit{\textbf{{\large ``Channel \hspace{1em} Abstractions"}}}}}}

%\newsavebox{\sstwolinebox}
%\sbox{\sstwolinebox}{\raisebox{-5pt}%
%{\parbox{7.4cm}{\textit{\textbf{{\large %
%Range-Bounded Types, Value Constructors, and Addressability}}}}}}
%\newcommand{\sstwoline}{\usebox{\sstwolinebox}}

%\definecolor{midRed}{rgb}{.5,.0,.25}
%\definecolor{logoRed}{rgb}{.3,0,0}



\newenvironment{displayquotexxx}{%
	%\begin{fcolorbox}{yellow!20!gray}{red!5}
	\begin{flushright}
	\begin{tcolorbox}[
		colback=gray!4,
		colframe=red!20!gray,
		width=.98\linewidth,
		boxrule=1pt,leftrule=3pt,
		rightrule=1pt,toprule=.5pt,arc=0pt,auto outer arc
		]
\begin{minipage}{21.5em}}{%
\end{minipage}\end{tcolorbox}
	\end{flushright}
}

\usepackage{changepage}

\newenvironment{dquote}{%
	%\begin{fcolorbox}{yellow!20!gray}{red!5}
	\begin{flushright}
	\vspace{1em}\begin{tcolorbox}[
		breakable, parbox=false, colback=gray!4,
		colframe=mRed!20!gray,
		width=.995\linewidth,
		boxrule=1pt,leftrule=3pt,
		rightrule=1pt,toprule=.5pt,arc=0pt,auto outer arc
		]
\begin{adjustwidth}{1pt}{-3pt}
}{%
\end{adjustwidth}
\end{tcolorbox}\vspace{1em}
	\end{flushright}
}



\newcounter{subs}[section]

\newcommand{\subss}[2]{%
\phantomsection \label{#1}
\addcontentsline{toc}{subsection}{#1}
\stepcounter{subs}
%\setcounterref{subsection}{\thesubs}
\refstepcounter{subsection}
\vspace*{3.25em}

\noindent{\LARGE\em\thesection.{\Large\thesubs}} {
\hspace{-2.5em}
			%\begin{minipage}{}
			\protect\raisebox{-5pt}{%
			\colorbox{logoCyan!50}{\hspace{2.1em}}}%
			\hspace{-5pt}{\protect\transparent{0.3}{\colorbox{logoBlue!80}{\protect\transparent{1}{%
						   \protect\raisebox{1pt}{\textit{{\large #2}}} }}}}
			%\end{minipage}

}
\vspace*{.35em} 
}

%\renewcommand{\subsection}[1]{\subss{#1}{#1}}


\newsavebox{\twolinebox}

\newcommand{\stwoline}[1]{%
\sbox{\twolinebox}{\raisebox{-3pt}%
{\parbox{7.4cm}{\linespread{1.25}\selectfont\raggedleft{\textit{\textbf{{\large #1}}}}}}}}

\newcommand{\twoline}{\usebox{\twolinebox}}

\newcommand{\subsectiontwoline}[1]{\stwoline{#1}\subss{#1}{\twoline}}

\newcommand{\spsubsectiontwolinerepl}[2]{\stwoline{#1}\subss{#2}{\twoline}}


\newcommand{\spsubsection}[1]{%
\subss{#1}{#1}
\vspace{1em}
}
\newcommand{\spsubsectiontwoline}[1]{%
\subsectiontwoline{#1}
\vspace{1em}
}



\usepackage{csquotes}

\usepackage{booktabs}

\usepackage{amssymb}

\usepackage{titlesec}

\usepackage{transparent}

\usepackage{setspace}

\usepackage{graphicx}

\newcommand{\sectionGraphics}{

\vspace{2em}\centerline{\includegraphics[scale=0.125]{logo-deco.png}}\vspace{-1em}}

\usepackage[flushmargin]{footmisc}

\setlength{\parindent}{30pt}

\usepackage[letterpaper, left=.47in,right=.47in,top=.94in,bottom=.94in]{geometry}

\let\xbibitem\bibitem
\renewcommand{\bibitem}[2]{\vspace{4pt}\xbibitem{#1}{#2}}

\newenvironment{frquote}{%
%\begin{fcolorbox}{yellow!20!gray}{red!5}
%\begin{flushright}
\begin{tcolorbox}[
	colback=orange!6,
	colframe=yellow!20!gray,
	width=.97\linewidth,
	arc=3mm, auto outer arc
	]
\begin{scriptsize}
\begin{minipage}{61em}
\begin{flushright}
\begin{minipage}{63em}}{%
\end{minipage}
\end{flushright}
\end{minipage}
\end{scriptsize}
\end{tcolorbox}
%\end{flushright}
}

%\colorlet{codegr}{black!80!blue}

\setlength{\columnsep}{8mm}

\usepackage{etoolbox}

\AtBeginEnvironment{thebibliography}{\linespread{1}\selectfont}

\usepackage{mathptmx}

\titleformat*{\subsection}{\small\bfseries}

\usepackage{wasysym}
\usepackage{textcomp}
\usepackage{amssymb}

\usepackage{microtype}

\DeclareMathAlphabet{\mathcal}{OMS}{cmsy}{m}{n}

\let\OldI\i

\newcommand{\secvspace}{\vspace{-0.3em}}
\newcommand{\asecvspace}{\vspace{-0.05em}}

\newcommand{\mdash}{---}
\newcommand{\q}[1]{``#1"}
\newcommand{\sq}[1]{`#1'}
\renewcommand{\i}[1]{\textit{#1}}

\newcommand{\nl}{

}

\newcommand{\T}[1]{\raisebox{-2pt}{\ensuremath{\mathcal{T}}}\textit{\tiny #1}}

\newcommand{\biburl}[1]{ {\fontfamily{gar}\selectfont{\textcolor[rgb]{.2,.6,0}%
{\scriptsize \textls*[-70]{\url{#1}}}}}}

\newcommand{\TSupT}{\ensuremath{{\T2}\makebox[4pt][r]{\raisebox{5pt}{{\scalebox{.6}{\T1}}}}}}

\let\OldFootnoteSize\footnotesize
\renewcommand{\footnotesize}{\scriptsize}
	
\newcommand{\Tnoindex}{\raisebox{-2pt}{\ensuremath{\mathcal{T}}}}

\newcommand{\typeAbove}{%
\raisebox{-1pt}{\rotatebox{90}{\begin{tiny}$\diagdown$\makebox[1pt][c]{$\diagup$}\end{tiny}}}}

\newcommand{\typeT}{\ensuremath{type\raisebox{.5pt}{\makebox[3pt][c]{-}}\mathcal{T}}}
\newcommand{\TValues}{\ensuremath{\mathcal{T}}-values} 

\newcommand{\emigres}{\'emigr\'es}

\newcommand{\SExpressions}{S-Expressions}
\newcommand{\SExpression}{S-Expression}

\newcommand{\cq}[1]{{\fontfamily{gar}\selectfont ``}#1{\fontfamily{gar}\selectfont "}}

\newcommand{\Retore}{Retor\'e}
\newcommand{\Aurelie}{Aur\'elie}
\newcommand{\Descles}{D\'escles}

\newcommand{\SG}{\ensuremath{\mathbb{SG}}}
\newcommand{\tLa}{\ensuremath{\mathbb{A}_\lambda}}

\newcommand{\lt}{\ensuremath{<}}
\newcommand{\ancestorLt}{\ensuremath{\lll}}
\newcommand{\aALTb}{\ensuremath{a {\lll} b}}
\newcommand{\aLTb}{\ensuremath{a {\lt} b}}
\newcommand{\aLTcLtb}{\ensuremath{a {\lt} b {\lt} c}}
\newcommand{\aDLTb}{\ensuremath{a {\lessdot} b}}

\newcommand{\aMath}{\ensuremath{a}}
\newcommand{\bMath}{\ensuremath{b}}
\newcommand{\cMath}{\ensuremath{c}}

\newcommand{\ala}{\`a la}

\newcommand{\br}{

}


\newcommand{\CppEleven}{{\Cpp}11}

\usepackage{graphicx}

\usepackage[breakable]{tcolorbox}

\newsavebox\lstbox



\newif\iffootnote
\let\Footnote\footnote
\renewcommand\footnote[1]{\begingroup\footnotetrue\Footnote{#1}\endgroup}



\newcommand{\AcronymText}[1]{{\iffootnote\begin{scriptsize}{\textsc{#1}}\end{scriptsize}%
\else\begin{OldFootnoteSize}{\textsc{#1}}\end{OldFootnoteSize}\fi}}

\newcommand{\TCP}{\AcronymText{TCP}}
\newcommand{\GPS}{\AcronymText{GPS}}
\newcommand{\CppTwenty}{\AcronymText{C++20}}

\newcommand{\AcronymTextInitialCap}[1]{{\iffootnote\begin{scriptsize}{\textsc{#1}}\end{scriptsize}%
		\else\begin{normalsize}{\textsc{#1}}\end{normalsize}\fi}}


\newcommand{\rel}[1]{\raisebox{0.25pt}{%
	{\iffootnote\begin{footnotesize}{\textsc{#1}}\end{footnotesize}%
		\else\begin{OldFootnoteSize}{\textsc{#1}}\end{OldFootnoteSize}\fi}}}

\newcommand{\vs}{

\vspace*{0.2em}}


\newcommand{\HTML}{\AcronymText{HTML}}

\newcommand{\UX}{\AcronymText{UX}}

\newcommand{\Csh}{\AcronymText{C\#}}

\newcommand{\PHP}{\AcronymText{PHP}}
\newcommand{\IDE}{\AcronymText{IDE}}

\newcommand{\AND}{\AcronymText{AND}}
\newcommand{\DX}{\AcronymText{DX}}

\newcommand{\RelaeGraph}{\AcronymText{RelaeGraph}}
\newcommand{\PDF}{\AcronymText{PDF}}

\newcommand{\NDP}{\AcronymText{NDP}}
\newcommand{\UML}{\AcronymText{UML}}
\newcommand{\GIT}{\AcronymText{GIT}}

\newcommand{\UI}{\AcronymText{UI}}
\newcommand{\QML}{\AcronymText{QML}}

\newcommand{\SPARQL}{\AcronymText{SPARQL}}
\newcommand{\OWL}{\AcronymText{OWL}}
\newcommand{\JSON}{\AcronymText{JSON}}

\newcommand{\const}{\AcronymText{const}}
\newcommand{\GeCODE}{\AcronymText{GeCODE}}

\newcommand{\Python}{\AcronymText{Python}}


\newcommand{\GUI}{\AcronymText{GUI}}
\newcommand{\BNF}{\AcronymText{BNF}}


\newcommand{\NLP}{\AcronymText{NLP}}

\newcommand{\Francois}{Fran\c{}cois}


\newcommand{\Gardenfors}{G\"ardenfors}


\newcommand{\Haskell}{\AcronymText{Haskell}}

\newcommand{\Cpp}{\AcronymText{C++}}
\newcommand{\RDF}{\AcronymText{RDF}}
\newcommand{\Java}{\AcronymText{Java}}
\newcommand{\IT}{\AcronymText{IT}}
\newcommand{\AI}{\AcronymText{AI}}

\newcommand{\SCA}{\AcronymText{SCA}}

\newcommand{\Lisp}{\AcronymText{Lisp}}

\newcommand{\Scheme}{\AcronymText{Scheme}}

\newcommand{\ThreeD}{\AcronymText{3D}}

\definecolor{BaseColor}{HTML}{8533FF}

\colorlet{ftcfore}{BaseColor!60!cyan}
\colorlet{ftcback}{BaseColor!40!cyan}

\newcommand{\ftc}[3]{
\vspace*{6mm}
\begin{tcolorbox}
[float=t, colframe=ftcfore!20!white,boxrule=0.5pt,arc=22pt,enhanced,
title={{\color{black}{\protect{#1}}}},label={#2}
toprule=0pt,bottomrule=1pt,
drop fuzzy shadow northeast={darkRed!40!ftcback},
      boxsep=3pt]\hspace{3em}\parbox{0.8\textwidth}\protect{#3}
\end{tcolorbox}      
\vspace*{-4mm}
}

\newcommand{\tc}[2]{
\vspace*{6mm}
\begin{tcolorbox}
[#1 colframe=darkRed!70!BaseColor,boxrule=0.5pt,arc=22pt,enhanced,
toprule=0pt,bottomrule=1pt,
drop fuzzy shadow northeast={darkRed!70!purple},
      boxsep=3pt]\hspace{3em}\parbox{0.8\textwidth}\protect{#2}
\end{tcolorbox}      
\vspace*{-4mm}
}

\tcbuselibrary{skins}
\usetikzlibrary{calc}
\usetikzlibrary{shadows}
\pgfdeclarelayer{background}
\pgfdeclarelayer{foreground}
\pgfsetlayers{background,main,foreground}



\usetikzlibrary{fit}
\usepackage{caption}

\renewcommand{\figurename}{Diagram}

\newcommand{\emblink}[2]{\href{#1}{#2}}

\newcommand{\tmphs}{\hypersetup{linkbordercolor=orange!50!red,linkcolor=black}}
\newcommand{\tmphscol}{\hypersetup{linkbordercolor=gray!40,linkcolor=black}}

\usepackage[object=vectorian]{pgfornament} 

\newcommand{\decoline}{%

\vspace{-2em}
{\color{darkRed!60!cyan}\noindent\hfil{\EnglischeLinie}\hfil}
\vspace{-2.25em}}


\newcommand{\whdecoline}{%

\vspace{-2em}
{\color{white}\noindent\hfil{\whEnglischeLinie}\hfil}
\vspace{-2.25em}}

\newcommand{\sectionline}[1]{%
  \noindent
  \begin{center}
  {\color{#1}
    \resizebox{0.5\linewidth}{1ex}
    {{%
    {\begin{tikzpicture}
    \node  (C) at (0,0) {};
    \node (D) at (9,0) {};
    \path (C) to [ornament=84] (D);
    \end{tikzpicture}}}}}%
    \end{center}
  }
  
\newcommand{\EnglischeLinie}{
\sectionline{darkRed!60!cyan}
}

\newcommand{\whEnglischeLinie}{
\sectionline{white}
}

\newcommand{\thinsectionline}[1]{%
	\noindent
	\begin{center}
		{\color{#1}
			\resizebox{.2\linewidth}{1.5ex}
			{{%
					{\begin{tikzpicture}
						\node  (C) at (0,0) {};
						\node (D) at (9,0) {};
						\path (C) to [ornament=84] (D);
						\end{tikzpicture}}}}}%
	\end{center}
}

\newcommand{\thindecoline}{\vspace*{-.15em}\thinsectionline{blGreen!10!orange}\vspace*{-.45em}}
\newcommand{\thindecolineadj}[2]{\vspace*{#1}\thinsectionline{blGreen!10!orange}\vspace*{#2}}


\newlength{\mplength}
\setlength{\mplength}{1.05\linewidth}

\newlength{\hslength}
\setlength{\hslength}{-1.45cm}

\newcommand{\spinc}[1]{\input{#1}}

\newsavebox{\tcsb}


\newcommand{\spinctc}[3]{\begin{lrbox}{\tcsb}\protect\input{#1}\end{lrbox} %\input{#1}
\begin{figure}
\tc{}{\protect\usebox\tcsb} 
\captionof{figure}{#2}
\label{#3}
\end{figure}
}

\usetikzlibrary{backgrounds} 
\usetikzlibrary{arrows}
\tikzset{>=triangle 90}

\newcommand{\gFunB}{\ensuremath{\mathfrak{g}}}
\newcommand{\fFunB}{\ensuremath{\mathfrak{f}}}

%\newcommand{\lstinlinebstyle}[1]{\Colorbox{blue!10}{%
%		{\fontfamily{Cabin-TLF}\fontseries{b}\fontsize{9pt}{16pt}\selectfont{{\textls[200]{#1}}}}}}

\newcommand{\lstinlinebstyle}[1]{\textls[200]{#1}}


\newenvironment{tikzgrid}{%
\begin{tikzpicture}[shorten >=1pt,node distance=2cm,on grid,auto,
framed,background rectangle/.style={double,ultra thick,draw=darkRed, bottom color=cyan!20, 
	top color=black!20, rounded corners}]%
}{\end{tikzpicture}}
 
\usepackage[font=small,labelfont=bf]{caption}

\definecolor{BaseColor}{HTML}{8533FF}

\colorlet{ftcfore}{BaseColor!60!cyan}
\colorlet{ftcback}{BaseColor!40!cyan}


\newcommand{\tcl}[1]{\begin{lrbox}{\lstbox}\begin{minipage}{\mplength}
\input{#1}
\end{minipage}\end{lrbox}
\tc{}{\hspace*{\hslength}\usebox\lstbox}
}

\newcommand{\itcl}[1]{
\input{#1}
}

\newcommand{\itclfig}[2]{
\begin{figure}\input{#1}
\label{#2}
\end{figure}
}

\newcommand{\ftcl}[2]{\begin{lrbox}{\lstbox}\begin{minipage}{\mplength}
\input{#2}
\end{minipage}\end{lrbox}
\tc{float=#1,}{\hspace*{\hslength}\usebox\lstbox}
}

\newcommand{\vsftcl}[2]{\begin{lrbox}{\lstbox}\begin{minipage}{\mplength}
\vsinput{#2}
\end{minipage}\end{lrbox}
\tc{float=#1,}{\hspace*{\hslength}\usebox\lstbox}
}

\newcommand{\vstcl}[1]{\begin{lrbox}{\lstbox}\begin{minipage}{\mplength}
%\vsinput{#1}
\end{minipage}\end{lrbox}  
%\tc{}{\hspace*{\hslength}\usebox\lstbox}
}

\definecolor{DarkRed}{rgb}{.2,.0,.1}

\definecolor{blGreen}{rgb}{.2,.7,.3}
\definecolor{darkRed}{rgb}{.2,.0,.1}

\definecolor{darkBlGreen}{rgb}{.1,.3,.2}

\definecolor{oldBlColor}{rgb}{.2,.7,.3}

\definecolor{blColor}{rgb}{.1,.3,.2}

\definecolor{elColor}{rgb}{.2,.1,0}
\definecolor{flColor}{rgb}{0.7,0.3,0.3}

\definecolor{logoOrange}{RGB}{108, 18, 30}
\definecolor{logoGreen}{RGB}{85, 153, 89}
\definecolor{logoPurple}{RGB}{200, 208, 30}

\definecolor{logoBlue}{RGB}{4, 2, 25}
\definecolor{logoPeach}{RGB}{255, 159, 102}
\definecolor{logoCyan}{RGB}{66, 206, 244}
\definecolor{logoRed}{rgb}{.3,0,0}

\definecolor{mRed}{rgb}{.5,.0,.25}



\definecolor{Bkg}{RGB}{250,245,252}
\newcommand{\leader}[2]{\hspace{#1}\colorbox{Bkg}{#2}}


\newcommand{\saying}[1]{\vspace{2ex}\noindent{%%
				\leader{2em}{\begin{minipage}{.93\textwidth}{\footnotesize #1}\end{minipage}}\vspace{2ex}}}

\newcommand{\sayingsrc}[1]{\vspace{0ex}\\\hspace{2pt} --- {\footnotesize #1}}

\renewcommand{\labelitemi}{{\raisebox{4pt}{{\footnotesize{$\bullet$}}}}}

\newcommand{\itemmark}{\raisebox{-4pt}{\rotatebox{90}{{\Large $\bracevert$}}}}

\usepackage{tikz}
\usetikzlibrary{positioning}
\usetikzlibrary{shapes,snakes}

\newcommand{\visavis}{vis-\`a-vis}

\newcommand{\tinyurl}[1]{{\raisebox{2pt}{{\scriptsize \url{#1}}}}} 

\usepackage[colorlinks=true]{hyperref}

\usepackage{eso-pic}

\colorlet{urlclr}{red!40!magenta!50!orange}

\hypersetup{
 urlcolor = urlclr,
 urlbordercolor = cyan!60!black,
 linkcolor = red!30!black,
 citecolor = orange!30!black,
 citebordercolor = yellow!30!black,
} 



\newcommand{\p}[1]{
	
	\vspace{.65em}
	#1	
}

\let\OldEnumerate\enumerate
  

\let\OldSection\section

\renewcommand\section[1]{
	\vspace{12pt}
	
	\scalebox{1.3}{\colorbox{logoPurple!50}{\hspace{1em}}}
					
	{\protect\transparent{0.5}{\colorbox{logoPeach}{%
			\begin{minipage}{\linewidth}
				    \vspace{.5em}
				
	\protect\transparent{1}{\OldSection{#1%
     }}
	    \end{minipage}}} }
    
    \vspace{-5em}
    	{\protect\transparent{0.3}{\colorbox{logoPeach}{%
    		\begin{minipage}{\linewidth}
    			\hspace{\linewidth}
    	\end{minipage}}} } 
    
    \vspace{5em}
    
	\vspace{-6pt}
}


\newcommand\tabltxt[1]{\footnotesize{#1}}
\renewcommand\section[1]{\OldSection{#1}}

\titleformat*{\subsection}{\Large\bfseries}

\let\OldSubsection\subsection
\renewcommand\subsection[1]{

	\vspace{12pt}
	
		\OldSubsection{% 	
			\hspace{-2.75em}
			\protect\raisebox{-5pt}{%
			\colorbox{logoCyan!50}{\hspace{2.1em}}}%
			\hspace{-5pt}{\protect\transparent{0.3}{\colorbox{logoBlue!80}{\protect\transparent{1}{%
						   \protect\raisebox{1pt}{\textit{{\large #1}}} }}}}
		}
	\vspace{-6pt}
}


\let\OLDthebibliography\thebibliography
\renewcommand\thebibliography[1]{
\let\section\OldSection
\setlength{\leftmargin}{-4pt}
\vspace{.1em}
\OLDthebibliography{#1}
\vspace{.7em}
\OldFootnoteSize 
\setlength{\parskip}{0pt}
\setlength{\itemsep}{1pt plus 0ex}
\raggedright
}

\makeatletter
\def\@biblabel#1{\hspace{-6pt}#1}
\makeatother

\newcommand{\bibtitle}[1]{{\small \textit{#1}}}
\newcommand{\intitle}[1]{{\hspace{3pt}\textls*[-80]{\texttt{\textit{#1}}}}\hspace{-1pt}}

\renewcommand{\i}[1]{\textit{#1}}

\usepackage{enumitem}
\setlist{leftmargin=1mm}



\newcommand{\itemtitle}[1]{{\color{green!10!red!40!black} \textls*[-80]{\texttt{#1}}}}

\newcommand{\acronymText}[1]{\AcronymText{#1}}
\newcommand{\acronymTextInitialCap}[1]{\AcronymTextInitialCap{#1}}
\newcommand{\acronymTextSS}[1]{\textsc{#1}}



\newcommand{\MIT}{\AcronymText{MIT}}

\newcommand{\nulltt}{\AcronymText{\texttt{null}}}

\newcommand{\boostUnits}{\codeText{\texttt{boost::units}}}

\newcommand{\Maybe}{\codeText{Maybe}}
\newcommand{\bind}{\codeText{bind}}
\newcommand{\return}{\codeText{return}}

\newcommand{\elseif}{\codeText{else if}}

\newcommand{\IDL}{\AcronymText{IDL}}


\newcommand{\DTD}{\AcronymText{DTD}}

\newcommand{\XML}{\AcronymText{XML}}
\newcommand{\API}{\AcronymText{API}}
\newcommand{\IR}{\AcronymText{IR}}
\newcommand{\Clang}{\AcronymText{Clang}}

\newcommand{\HCI}{\AcronymText{HCI}}
\newcommand{\HTTP}{\AcronymText{HTTP}}
\newcommand{\CSS}{\AcronymText{CSS}}

\newcommand{\STL}{\AcronymText{STL}}

\newcommand{\Qt}{\AcronymText{Qt}}
\newcommand{\NL}{\AcronymText{NL}}

\newcommand{\vsinput}[1]{\vspace{1em}\input{#1}}

\newcommand{\qi}[1]{\q{\textit{#1}}}

\newcommand{\CSharp}{\AcronymText{C\#}}
\newcommand{\IoT}{\AcronymText{IoT}}

\usetikzlibrary{decorations.pathmorphing}

\usepackage{listings}

\renewcommand{\lstlistingname}{Sample}

%%%%%%%%
\newcommand{\OO}{\AcronymText{OO}}
\newcommand{\SQL}{\AcronymText{SQL}}
\newcommand{\QtCpp}{\AcronymText{Qt/C++}}
\newcommand{\JavaScript}{\AcronymText{JavaScript}}

\newcommand{\ECL}{\AcronymText{ECL}}

\newcommand{\ovn}[1]{\color{yellow}{{\textbf{#1}}}}

\newcommand{\sovn}[1]{\color{yellow}{{\textbf{#1}}}}

\newcommand{\dovn}[3]{\draw[draw=blue,fill=DarkRed] (#1,#2) circle[radius=3mm];
\node at (#1,#2){\ovn{#3}}
}

\newcommand{\sdovn}[3]{\draw[draw=blue,fill=DarkRed] (#1,#2) circle[radius=1.4mm];
\node (char) at (#1,#2) {\sovn{#3}}
}

\newcommand{\sdovnz}[3]{\raisebox{-1pt}{{\draw[draw=blue,fill=DarkRed] (#1,#2) circle[radius=1.4mm];
\node (char) at (#1,#2) {\sovn{#3}}
}}}

\newcommand*\circledx[1]{\tikz[baseline=(char.base), inner sep=0]{
            \sdovn{0}{0}{#1};}}

\newcommand{\circled}[1]{\raisebox{-.5pt}{\circledx{#1}}}

\newcommand{\circledup}[1]{\circledx{#1}}

\newcommand{\circledd}[1]{\raisebox{-2pt}{\circledx{#1}}}
			
\newcommand{\OneOverlay}{\circled{1}}
\newcommand{\TwoOverlay}{\circled{2}}
\newcommand{\ThreeOverlay}{\circled{3}}
\newcommand{\FourOverlay}{\circled{4}}
\newcommand{\FiveOverlay}{\circled{5}}
\newcommand{\SixOverlay}{\circled{6}}

\newcommand{\OneOverlayu}{\circledup{1}}
\newcommand{\TwoOverlayu}{\circledup{2}}
\newcommand{\ThreeOverlayu}{\circledup{3}}
\newcommand{\FourOverlayu}{\circledup{4}}
\newcommand{\FiveOverlayu}{\circledup{5}}
\newcommand{\SixOverlayu}{\circledup{6}}

\newcommand{\FourOverlayd}{\circledd{4}}

\newcommand{\true}{\codeText{true}}

\colorlet{codegr}{black!80!blue}

\newcommand{\FnDoc}{\codeText{Fn\_Doc}}

\newcommand{\KCMEnv}{\codeText{KCM\_Env}}

\newcommand{\envv}{\codeText{envv}}

\newcommand{\kprom}{\codeText{kcm\_promote\_type\_binding\_to%
\_statement\_via\_type\_de\-fault\_literal\_\_let}}
\newcommand{\kcmde}{\codeText{kcm\_direct\_eval}}

\newcommand{\dbleq}{\codeText{==}}
\newcommand{\Sdbleq}{\codeText{$\backslash$==}}
\newcommand{\Seq}{\codeText{$\backslash $=}}

\usepackage[utf8]{inputenc}
 
\usepackage{listings}
\usepackage{color}
 
\definecolor{codegreen}{rgb}{0,0.6,0}
\definecolor{codegray}{rgb}{0.5,0.5,0.5}
\definecolor{codepurple}{rgb}{0.58,0,0.82}
\definecolor{backcolour}{rgb}{0.95,0.95,0.92}
 
\usetikzlibrary{positioning}

\usepackage{listings}
\lstset{%
  frame            = tb,    % draw frame at top and bottom of code block
  tabsize          = 1,     % tab space width
  numbers          = left,  % display line numbers on the left
  framesep         = 3pt,   % expand outward
  framerule        = 0.4pt, % expand outward 
  commentstyle     = \color{Green},      % comment color
  keywordstyle     = \color{blue},       % keyword color
  stringstyle      = \color{DarkRed},    % string color
  backgroundcolor  = \color{gray!20}, % backgroundcolor color
  showstringspaces = false,              % do not mark spaces in strings
  upquote          = true
}

\newcommand{\yeqfxz}{\codeText{y=f(x,z)}}

\newcommand{\try}{\codeText{try}}
\renewcommand{\DH}{\AcronymText{DH}}
\newcommand{\CH}{\AcronymText{CH}}

%\let{\oldeta}{\eta}
%\renewcommand{\eta}{\ensuremath{eta}}


\newcommand{\catch}{\codeText{catch}}


\newcommand{\vecrgbt}{\AcronymText{RGBT}}
\newcommand{\vecrgb}{\AcronymText{RGB}}

\newcommand{\TTL}{\AcronymText{TTL}}

\newcommand{\URL}{\AcronymText{URL}}
\newcommand{\RRF}{\AcronymText{RRF}}
\newcommand{\SPO}{\angla{\AcronymText{Subject}, %
\AcronymText{Predicate}, \AcronymText{Object}}}

\newcommand{\Turtle}{\AcronymText{Turtle}}

\newcommand{\CLang}{\AcronymText{C}}

 
\newcommand{\intxeqninety}{\codeText{int x = 90}}

\newcommand{\tyFile}{\codeText{file}}
\newcommand{\idrisText}[1]{\codeText{#1}}
\newcommand{\openFn}{\codeText{open}}

\newcommand{\ML}{\acronymText{ML}}

\newcommand{\largeRDF}{RDF}

\newcommand\Small{\fontsize{8}{10}\selectfont}
\newcommand*\LSTfont{%
  \Small\ttfamily\SetTracking{encoding=*}{-50}\lsstyle}

	
\newcommand{\NathLivedTTL}{\lstset{ basicstyle=\LSTfont, columns=fullflexible, xleftmargin=5mm, framexleftmargin=5mm, numbers=left, stepnumber=1, breaklines=true, breakatwhitespace=false, numberstyle=\footnotesize, numbersep=5pt, 
tabsize=2, frame=lines, captionpos=b, caption={Turtle Formats}}
    %\lstinputlisting{NathLivedTTL.n3}
    }		

\newcommand{\TyS}{\codeTextr{$\mathbb{T}$}}

\newcommand{\TXLTyS}{$\codeTextr{\mathfrak{L}_\codeTextr{\mathbb{T}}}$}

\newcommand{\TXLTySChi}{$\codeTextr{\mathfrak{L}_\codeTextr{\mathbb{T}}\chiussr}$}


\newcommand{\ageFF}{\codeTexto{${\lceil}45{\rceil}$}}
\newcommand{\Nath}{\codeTexto{${\lceil}$Nathaniel${\rceil}$}}

\newcommand{\nodeNOne}{$N_1$}
\newcommand{\nodeNTwo}{$N_2$}

\newcommand{\NThree}{\AcronymText{N3}}

\newcommand{\ceila}[1]{${\lceil}$#1${\rceil}$}
\newcommand{\angla}[1]{${\langle}$#1${\rangle}$}

\newcommand{\NathFF}{\codeTexto{\angla{\ceila{Nathaniel}, \ceila{46}}}}
\newcommand{\NathFFBD}{\codeTexto{\angla{\ceila{Nathaniel}, \ceila{46}, %
\ceila{Brooklyn}, \ceila{Democrat}}}}

\newcommand{\BrookDem}{\codeTexto{\angla{\ceila{Brooklyn}, \ceila{Democrat}}}}

\newcommand{\nameAge}{\codeTexto{${\langle}$\codeText{name}, \codeText{age}${\rangle}$}}

\newcommand{\suigeneris}{\i{sui generis}}
 
\newcommand{\struct}{\codeText{struct}}
\newcommand{\CStructs}{\codeText{struct}s}
\newcommand{\CStructsArrays}{\codeText{struct}s/arrays}

\newcommand{\CStrucst}{C\codeText{struct}}

\newcommand{\throw}{\codeText{throw}}
\newcommand{\exception}{\codeText{exception}}

\newcommand{\float}{\codeText{float}}

\newcommand{\chiuss}{\raisebox{-1pt}{$^\chiu$}}
\newcommand{\chiussr}{\raisebox{-1pt}{$^\chiur$}}

\newcommand{\chiur}{\codeTextr{\ensuremath{\chi}}}
\newcommand{\chiu}{\codeText{\ensuremath{\chi}}}

\newcommand{\TySChi}{\TyS\chiussr}


\newcommand{\lCh}{\chsnt{lam}}
\newcommand{\rCh}{\chsnt{ret}}
\newcommand{\xCh}{\chsnt{exc}}

\newcommand{\lrCh}[1]{\chsnt{lam, ret, #1}}

\newcommand{\lrChwow}[1]{\chsnt{lam!, ret!, #1}}

\newcommand{\lsrxCh}{\chsnt{lam, sig, ret, exc}}

\newcommand{\lrChblank}{\chsnt{lam, ret}}

\newcommand{\rChSize}{\chsntsz{ret}}
\newcommand{\chChSize}{\chsntsz{ch}}

\newcommand{\rNoMixx}{\rCh{\chcolor{$\nshortparallel$}}\xCh}
\newcommand{\rOneOrx}{\rCh{\chcolor{$\asymp$\xCh}}}

\newcommand{\rChSizeleOne}{\codeText{\rChSize{} $\leq$ 1}}
\newcommand{\rChSizeLe}{\codeText{rChSizele}}
\newcommand{\Ch}{\chsnt{ch}}

\newcommand{\lrxSimple}{\lrChwow{exc?}}

\newcommand{\lrxDetailed}{\makebox{%
\chsnt{lam!*},\chnt{ret!\%},\chnt{exc?\%}}}

\newcommand{\excl}{\chcolor{\ensuremath{\asymp}}}

\newcommand{\lrxTotal}{\makebox{%
\lrCh{exc} {\colonblg} %
\chnt{lam!*},\chnt{ret!\%},\chnt{exc?\%}%
{ }{\colonblg}{\colonblg} \chnt{ret}{\excl}\chnt{exc}}}

\newcommand{\lr}{\makebox{%
\lrChblank{}{ }{\colonblg} \chnt{lam!*},\chnt{ret!\%}}}

\newcommand{\lsrx}{\makebox{\makebox{\lsrxCh}%
{ }{\colonblg} \chnt{lam!*},\chnt{sig?\%},\chnt{ret!\%},\chnt{exc?\%}}}
\newcommand{\sCh}{\chsnt{sig}}


\newcommand{\lrx}{\makebox{%
\lrCh{exc}}}

\let\OldLambda\lambda

\renewcommand{\lambda}{\codeText{lambda}}
\renewcommand{\return}{\codeText{return}}

\newcommand{\lambdaPLUSreturn}{{\lambda}{\codeTextr{+}}{\return}}


\newcommand{\lxr}{\codeText{lxr}}

\newcommand{\fnote}[1]{\codeText{#1}}

\newcommand{\codeinclude}{\codeText{\#include}}
\newcommand{\codeconstruct}{\codeText{construct}}

\newcommand{\USH}{\acronymText{USH}}

\newcommand{\fxy}{\codeText{\makebox{f(x, y)}}}
\newcommand{\fSym}{\codeText{f}}

\newcommand{\fofx}{\codeText{\makebox{f(x)}}}
\newcommand{\fofy}{\codeText{\makebox{f(y)}}}

\newcommand{\fFuns}{\codeText{f}}
\newcommand{\hFun}{\codeText{h}}
\newcommand{\hfx}{\codeText{h(f(x))}}
\newcommand{\fxdoth}{\codeText{f(x).h()}}

\newcommand{\typesH}{\codeText{types.h}}

\newcommand{\QDataStream}{\codeText{QDataStream}}

\newcommand{\CATSet}{$\mathbb{S}$}

\newcommand{\DSL}{\AcronymText{DSL}}
\newcommand{\TXL}{\AcronymText{TXL}}
\newcommand{\CXL}{\AcronymText{CXL}}

\newcommand{\IntZToOH}{\codeText{int\ranged{0,100}}}

\newcommand{\OH}{\codeText{100}}

\newcommand{\codebreak}{\codeText{break}}
\newcommand{\codecontinue}{\codeText{continue}}

\newcommand{\TVOneToVTwo}{\codeText{T\ranged{$V_1$,$V_2$}}}

\newcommand{\RangeLTEVal}{\codeText{ranged\_lte}}
\newcommand{\RangeLTEOHxeqOHone}{\codeText{ranged\_lte{$<$100$>$} x = 101}}
\newcommand{\RangeLTEValV}{\codeText{\codeText{ranged\_lte{$<$V$>$}}}}

\newcommand{\MIpair}{\codeText{mi\_pair}} 
\newcommand{\fmipair}{\codeText{int f(mi\_pair pr)}}

\newcommand{\fileTy}{\codeText{file}} 
\newcommand{\pairOfLists}{\codeText{pair$<$list$<$...$>>$}}


\newcommand{\VOne}{\codeText{$V_1$}}
\newcommand{\VTwo}{\codeText{$V_2$}}

\newcommand{\TType}{\codeText{T}}

\newcommand{\crVOverRTwo}{\codeTextr{$\sqrt{3}\sqrt[3]{V}$}}

\newcommand{\volSphCube}{\codeTextr{$\frac{4}{3}\sqrt{27}\pi V$}}

\newcommand{\rRad}{\codeTextr{$R$}}
\newcommand{\vVol}{\codeTextr{$V$}}
\newcommand{\piSym}{\codeTextr{$\pi$}}


\newcommand{\TMyList}{\codeText{MyList$<$T$>$}}

\newcommand{\MyList}{\codeText{MyList}}

\newcommand{\templateTMyList}{\codeText{template$<$T$>$MyList}}
\newcommand{\MyListInt}{\codeText{MyList$<$int$>$}}

\newcommand{\listsize}{\codeText{list.size()}}
\newcommand{\sizelist}{\codeText{size(list)}}

\newcommand{\SLE}{\AcronymText{SLE}}

\newcommand{\tCat}{\TyS}
\newcommand{\DBH}{\AcronymText{DBH}}

\newcommand{\cCat}{$\mathbb{C}$}
\newcommand{\eOneToeTwo}{$e_1 \rightarrow e_2$}
\newcommand{\eOne}{$e_1$}
\newcommand{\eTwo}{$e_2$}

\newcommand{\VolSphere}{\codeTextr{$\frac{{4\pi R^3 }}{3}$}}

\newcommand{\livedinup}{\raisebox{6pt}{\colorbox{blGreen!20}{lived-in}}}

\newcommand{\rdfrestup}{\raisebox{4pt}{\colorbox{blGreen!20}{rdf:rest}}}


\newcommand{\CArray}{\AcronymText{C} array}
\newcommand{\CStruct}{\AcronymText{C} \codeText{struct}}


\newcommand{\Chi}{$\codeTextr{\chi}$}

\newcommand{\NTrips}{\AcronymText{N-Triples}}
\newcommand{\NTrip}{\AcronymText{N-Triple}}

\renewcommand{\GeCODE}{\AcronymText{GeCODE}}
\newcommand{\ITK}{\AcronymText{ITK}}

\newcommand{\SSI}{\AcronymText{SSI}}
\newcommand{\GSI}{\AcronymText{GSI}}

\newcommand{\ISBN}{\AcronymText{ISBN}}
\newcommand{\GIS}{\AcronymText{GIS}}


\newcommand{\CodeMinted}[1]{{\color{codegr}{{%
\fontfamily{lmss}\fontseries{b}\selectfont{#1}}}}}

\newcommand{\CodeMintedo}[1]{{\color{orange!40!black}{{%
\fontfamily{lmss}\fontseries{b}\selectfont{#1}}}}}

\newcommand{\CodeMintedr}[1]{{\color{red!40!black}{{%
\fontfamily{lmss}\fontseries{b}\selectfont{#1}}}}}

\newcommand{\codeText}[1]{\CodeMinted{#1}}
\newcommand{\codeTexto}[1]{\CodeMintedo{#1}}
\newcommand{\codeTextr}[1]{\CodeMintedr{#1}}

\newcommand{\FactPP}{\AcronymText{Fact++}}

\newcommand{\chname}[1]{\AcronymText{\textbf{#1}}}

\newcommand{\chsym}{\raisebox{3pt}{\rotatebox{-45}{$\Arrowvert$}}\hspace{-5pt}%
\raisebox{1pt}{\rotatebox{-45}{\tiny{$\gg$}}}\hspace{-1pt}}  %\searrow

\newcommand{\chsymt}{\raisebox{3pt}{\rotatebox{-45}{$\Arrowvert$}}\hspace{-5pt}%
\raisebox{1pt}{\rotatebox{-45}{\tiny{$\gg$}}}\hspace{-1pt}}  %\searrow \Rightarrow

\newcommand{\chsnt}[1]{{\color{blGreen!40!black}{\chsymt\chname{#1}}}}
\newcommand{\chnt}[1]{{\color{blGreen!40!black}{\chname{#1}}}}

\newcommand{\colonblg}{{\color{blGreen!40!black}{:}}}

\newcommand{\chcolor}[1]{{\color{blGreen!40!black}{#1}}}

\newcommand{\chsntsz}[1]{{\color{blGreen!40!black}{\makebox{\#\chsymt\chname{#1}}}}}


\newcommand{\chsn}[1]{\chsym\chname{#1.}}

\newcommand{\lXY}{\codeTexto{${\OldLambda}x.{\OldLambda}y$}}


\newcommand{\CHlXY}{\chcolor{\chsn{lam}$xy$}}



\newcommand{\Tvar}{\codeText{T}}
\newcommand{\TrRan}{\codeText{T\ranged{r}}}
\newcommand{\rRan}{\codeText{\ranged{r}}}

\newcommand{\jFunction}{\ensuremath{j}-function}

\newcommand{\xVal}{\codeText{x}}
\newcommand{\xeqyplusz}{\codeText{x $=$ y $+$ z}}

\newcommand{\fOfG}{\codeText{$f{\circ}g$}}
\newcommand{\Ofop}{\codeText{$\circ$}}

\newcommand{\inc}{\codeText{inc}}

\newcommand{\zeroNum}{\codeText{$0$}}


\newcommand{\enumleft}{\codeText{$\langle$}}
\newcommand{\enumright}{\codeText{$\rangle$}}

\newcommand{\closed}{\codeText{closed}}
\newcommand{\open}{\codeText{open}}
\newcommand{\nonexistent}{\codeText{nonexistent}}

\newcommand{\codeDH}{\codeText{code}-{\DH}}


\newcommand{\fOfg}{\codeText{$f{\circ}g$}}  


\newcommand{\cfFun}{\codeText{Cf}}
\newcommand{\cf}{\codeText{Cf}}
\newcommand{\zToOH}{\codeText{\ranged{0, 100}}}
\newcommand{\C}{\codeText{C}}
\newcommand{\yplusz}{\codeText{y $+$ z}}
\newcommand{\iVal}{\codeText{$i$}}
\renewcommand{\le}{\codeText{$\leq$}}

\renewcommand{\int}{\codeText{int}}

\newcommand{\tTy}{\codeText{T}}

\newcommand{\rrRanOfTVV}{\codeText{ranged$<$T, t1, t2$>$}}

\newcommand{\nVal}{$n$}

\newcommand{\zTon}{\codeText{\ranged{0, n}}} 

\newcommand{\fOneTwoxeq}{\codeText{$f_1(x)=f_2(x)$}}

\newcommand{\fOne}{\codeText{$f_1$}}
\newcommand{\fTwo}{\codeText{$f_2$}}

\newcommand{\xVar}{\codeText{$x$}}

\newcommand{\fx}{\codeText{f(x)}}

\newcommand{\tOne}{\codeText{$\mathcal{T}_1$}}
\newcommand{\tTwo}{\codeText{$\mathcal{T}_2$}}

\newcommand{\gSym}{\codeText{g}}
\newcommand{\gradeFn}{\codeText{grade}}
\newcommand{\iSym}{\codeText{i}}

\newcommand{\xninetyyonehgrade}{\codeText{x $=$ 90; y $=$ 100; g $=$ grade(x, y)}}

\newcommand{\tOneTotTwo}{\codeText{$\mathcal{T}_1 \Rightarrow \mathcal{T}_2$}}

\newcommand{\zeroToOH}{\codeText{\ranged{0,100}}}

\newcommand{\rrsb}[1]{\raisebox{5pt}{#1}}
\newcommand{\ranged}[1]{\codeText{\rrsb{{\tiny{$\lgroup$}}}#1\rrsb{{\tiny{$\rgroup$}}}}}

\newcommand{\ftytwoh}{\codeTextr{{$\big[40-200\big]$}}}

\newcommand{\ZeroToOneHundred}{\codeText{\ranged{0,100}}}

\newcommand{\rRanOfT}{\codeText{ranged$<$T$>$}}


\newcommand{\xSym}{\codeText{x}}
\newcommand{\ySym}{\codeText{y}}
\newcommand{\zSym}{\codeText{z}}
\newcommand{\fFun}{\codeText{f}}

%%% SLE
\newcommand{\addressOf}{\codeText{address-of}}
\newcommand{\fofg}{\codeText{$f{\circ}g$}}

\newcommand{\smsp}{\hspace{2pt}}

\newcommand{\tOnetoTwotoThree}{\codeText{$\mathcal{T}_1$\smsp%
$\rightarrow$\smsp$\mathcal{T}_2$\smsp$\rightarrow$\smsp$\mathcal{T}_3$}}
 
\newcommand{\tOnetoTwoTOThree}{\codeText{$\mathcal{T}_1$\smsp%
$\rightarrow$\smsp\smsp($\mathcal{T}_2$\smsp$\rightarrow$\smsp$\mathcal{T}_3$)}} 
 
\newcommand{\pimath}{\ensuremath{\pi}}

\newcommand{\SCO}{\AcronymText{SCO}}

\newcommand{\lambdaxfgx}{\codeTextr{$\OldLambda{}x.fgx$}} 
\newcommand{\lambdaxfx}{\codeTextr{$\OldLambda{}x.fx$}} 

\newcommand{\gFun}{\codeText{g}}  

\newcommand{\nodex}{\codeText{$x$}}
\newcommand{\nodef}{\codeText{$f$}}

\newcommand{\sqqq}{{\fontfamily{lmtt}\selectfont{"}}}
\newcommand{\operatorqq}{\codeText{operator\sqqq\sqqq}}

\newcommand{\carOne}{\codeTextr{$\mathfrak{C}_1$}}

\newcommand{\carTwo}{\codeTextr{$\mathfrak{C}_2$}} 

\newcommand{\carOnetoTwos}{\carOne \codeTextr{$\twoheadrightarrow$} \carTwo}
\newcommand{\carOnetoTwof}{\carOne \codeTextr{$\rightarrowtail$} \carTwo}


\newcommand{\intieqzero}{\codeText{int i = 0}}

\newcommand{\intthrtwo}{\codeText{int32}}

\newcommand{\aeqb}{\codeText{a = b}}

\newcommand{\aceqb}{\codeText{a := b}}

\newcommand{\thisc}{\codeText{this}}

\newcommand{\lambdaCalculus}{$\OldLambda$-Calculus}

\newcommand{\lambdas}{$\OldLambda$s}


\newcommand{\fntoch}{\codeTextr{$f_n$ $\rightarrow$ $\chi$}}

\newcommand{\sigmac}{\codeTextr{sigma}}
\newcommand{\sigmach}{\codeTextr{sigma}} 

\newcommand{\sigmaCalculi}{$\varsigma$-calculi}
\newcommand{\sigmaCalculii}{$\varsigma$-calculi}
\newcommand{\sigmaCalculus}{$\varsigma$-calculus}

\newcommand{\returnc}{\codeTextr{return}}

\newcommand{\lambdach}{\codeTextr{lambda}}

\newcommand{\catchddd}{\codeTextr{catch(...)}}

\newcommand{\sectsym}{\S}

\newcommand{\error}{\codeTextr{exception}}
\newcommand{\errorc}{\codeTextr{exception}}
\newcommand{\errorrc}{\codeTextr{exception}}
\newcommand{\returnrc}{\codeTextr{return}}

\newcommand{\returnch}{\codeTextr{return}}
\newcommand{\fgroundch}{\codeTextr{fground}}


\newcommand{\lambdac}{\codeTextr{lambda}}
\newcommand{\capturec}{\codeTextr{capture}}
\newcommand{\capturech}{\codeTextr{capture}}

\newcommand{\errorch}{\codeTextr{exception}}
\newcommand{\breakch}{\codeTextr{break}}
 
\renewcommand{\break}{\codeText{break}}
\newcommand{\breakct}{\codeText{break}}

\newcommand{\yeqfx}{\codeText{y = f(x)}} 
\newcommand{\frcchoprctocy}{\codeText{frcchoprctocy}} 

\newcommand{\frc}{\codeText{frc}} 
\newcommand{\fFn}{\codeText{fFn}} 

\newcommand{\bindFn}{\codeText{bindFn}} 

\newcommand{\lookupch}{\codeTextr{lookup}} 


\newcommand{\stdFuture}{\codeText{std::future}} 
 
 
\newcommand{\doH}{\codeText{do}}

\newcommand{\ifthenelse}{\codeText{if...then...else}}

\newcommand{\tyListOfInt}{\codeText{list$<$int$>$}}
\newcommand{\tyint}{\codeText{int}}

\let\OldGamma\Gamma

\renewcommand{\Gamma}{\codeTextr{$\OldGamma$}}
\newcommand{\gammaOne}{\codeTextr{$\OldGamma_1$}}
\newcommand{\gamaTwo}{\codeTextr{$\OldGamma_2$}}
\newcommand{\gammaTwo}{\codeTextr{$\OldGamma_2$}}

\newcommand{\pfx}{\codeText{(*f)(x)}}

\newcommand{\oprctoc}{\codeText{oprctoc}}
\newcommand{\oprctoch}{\codeText{oprctoch}}
\newcommand{\choprctoc}{\codeText{choprctoc}}

\newcommand{\carOnetoTwofOP}{\codeText{carOnetoTwofOP}}

\newcommand{\thisob}{\codeText{this}}
\newcommand{\selfob}{\codeText{self}}

\newcommand{\typetoch}{\codeTextr{%
\raisebox{2pt}{$\leftharpoonup$}\hspace{-.9em}\raisebox{1pt}{$\leftharpoondown$}}}

\newcommand{\nodeT}{\codeTextr{T$_{Node}$}}
\newcommand{\predTY}{\codeTextr{T$_{Pred}$}}
\newcommand{\predTYOp}{\codeTextr{\lambdach. \typetoch T$_{Pred}$ %
\returnch. \typetoch T$_{Pred}$}}

\newcommand{\predTYsig}{\sigmach\codeTextr{\%}\returnch\codeTextr{\%}}

\newcommand{\QFile}{\codeText{QFile}}
  

\newcommand{\xCommaY}{\codeText{$x,y$}}  
\newcommand{\xlty}{\codeText{x $<$ y}}   
\newcommand{\unsignedint}{\codeText{unsigned int}}
 
\newcommand{\subjectPredicateObject}{\codeTexto{Subject Predicate Object}}
\newcommand{\rdfType}{\codeTexto{rdf:type}}

\newcommand{\parentOf}{\codeTexto{parentOf}}
\newcommand{\childOf}{\codeTexto{childOf}}
\newcommand{\inverseOp}{\codeTexto{inverseOp}}

\newcommand{\capabilitych}{\codeTextr{capability}} 
 
\newcommand{\addFun}{\codeText{add}} 

\newcommand{\sortFun}{\codeText{sort}}
\newcommand{\sortfn}{\codeText{sort}} 

\newcommand{\incimpl}{\codeText{int inc(int x)\{return add(x,1)\}}} 
 
\newcommand{\addOne}{$\langle$\codeText{\&add, 1}$\rangle$}

\newcommand{\Cfr}{$\langle$\codeText{\&Cf, r$_1$, r$_2$}$\rangle$}


\newcommand{\this}{\codeText{this}} 
 
\newcommand{\fsVal}{\codeText{fs}} 
\newcommand{\static}{\codeText{static}}
 
\newcommand{\vVal}{\codeText{V}}  
\newcommand{\lte}{\codeText{$\leq$}}
\newcommand{\lteVal}{\codeText{$\leq{ }$V}}

\newcommand{\intXeqOPointFive}{\codeText{int x = $0.5$}} 

\newcommand{\RangeGTVal}{\codeText{ranged\_gt}}
\newcommand{\RangeGTValx}{\codeText{ranged\_gt{$<$x$>$}}}

\newcommand{\cFun}{\codeText{C}}

\newcommand{\st}{\codeTextr{s.t}} 

\newcommand{\prVal}{\codeText{pr}}

\newcommand{\parach}{\codeTextr{para}}

\newcommand{\fxypreqMIpairxy}{\codeText{f(int x, int y, mi\_pair pr = mi\_pair(x, y))}}
\newcommand{\fxypreqMIpairZeroOne}{\codeText{f(x, y, mi\_pair(0, 1))}}
 
 
%\newcommand{\fxypreqMIpairxy}{\codeText{f int x, int y, mi\_pair pr = mi\_pair x, y}}
%\newcommand{\fxypreqMIpairZeroOne}{\codeText{def}}
 
 
 
\newcommand{\yVal}{\codeText{y}}
\newcommand{\xSymbol}{\codeText{x}}
 
\newcommand{\subjectnd}{\codeTexto{SUBJECT}}
\newcommand{\predicatend}{\codeTexto{PREDICATE}}
\newcommand{\objectnd}{\codeTexto{OBJECT}}
 
\newcommand{\doNotation}{\codeText{do}-notation}

\newcommand{\Onef}{\codeText{$f_1$}}

\newcommand{\Twof}{\codeText{$f_2$}}

\newcommand{\CLng}{\AcronymText{C}}

\newcommand{\qOneTwoThree}{\q{\codeText{123}}} 
\newcommand{\oneTwoThree}{\codeText{123}}
\newcommand{\ascii}{\AcronymText{ASCII}}
 
 
\newcommand{\opTransfer}{\codeTextr{$\looparrowright$}}
\newcommand{\carrOne}{\codeTextr{$\mathfrak{C}_1$}}
\newcommand{\carrTwo}{\codeTextr{$\mathfrak{C}_2$}}

\newcommand{\catchexce}{\codeText{catch(Exception e)}}

\newcommand{\carrOneOpTransferTwo}{\carrOne{}{\opTransfer}\carrTwo}
\newcommand{\carrOneOpTransferTwolambda}{\carrOne{}{\opTransfer}%
{\textsuperscript{\hspace{-.8em}{\lambdach}}}\carrTwo}


\newcommand{\carrOneOpTransferTworeturn}{\carrOne{}{\opTransfer}%
{\textsuperscript{\hspace{-.8em}{\returnch}}}\carrTwo}

\newcommand{\carrTwoOpTransferOnereturn}{\carrTwo{}{\opTransfer}%
{\textsuperscript{\hspace{-.8em}{\returnch}}}\carrOne}

\newcommand{\exceptionch}{\codeTextr{exception}}

\newcommand{\chanOne}{\codeTextr{$\chi_1$}} 
\newcommand{\chanTwo}{\codeTextr{$\chi_2$}} 

\newcommand{\chanOneOpTransferTwo}{\chanOne{}{\opTransfer}\chanTwo}

\newcommand{\objfx}{\codeText{obj.f(x)}}
\newcommand{\obj}{\codeText{obj}}

%\newcommand{\chanyOne}{\codeTextr{$\chi_1$}} 
%\newcommand{\changTwo}{\codeTextr{$\chi_2$}} 
%\newcommand{\chanyOneOpTransferTwo}{\chanOne{}{\opTransfer}\chanTwo}



\newcommand{\gOpTransferOneOneF}{\codeText{g:}\codeTextr{return$_1$}{\opTransfer}%
\codeText{f:}\codeTextr{lambda$_1$}}
 
\newcommand{\fDotOfg}{\codeText{f$\circ$g}}
\newcommand{\fDotOfGX}{\codeText{(f.g)(x)}}
\newcommand{\fOfGx}{\codeText{f(g(x))}}

\newcommand{\fgx}{\codeText{f(g(x))}}

\newcommand{\funG}{\codeText{g}}
\newcommand{\funF}{\codeText{f}} 

\newcommand{\tyOne}{\codeText{$\mathcal{T}_1$}}
\newcommand{\tyTwo}{\codeText{$\mathcal{T}_2$}}

\newcommand{\tyOneTotyTwo}{\tyOne \codeText{$\rightarrow$} \tyTwo}

\newcommand{\chK}{\codeText{$\mathcal{K}$}}

\newcommand{\tOneTimesTTwo}{\codeText{$\mathcal{T}_1$ $\times$ $\mathcal{T}_2$}}
\newcommand{\tOneTimesTOne}{\codeText{$\mathcal{T}_1$ $\times$ $\mathcal{T}_1$}}
\newcommand{\tyOneTimesTyTwo}{\codeText{$\mathcal{T}_1$ $\times$ $\mathcal{T}_2$}}

\newcommand{\tOneTimesTTwoToTOneOntoTTwo}{\codeText{$\mathcal{T}_1$ $\times$ $\mathcal{T}_2$ %
$\Longrightarrow$ $\mathcal{T}_1$ $\rightarrow$ $\mathcal{T}_2$}}

\newcommand{\unitTy}{\codeText{Unit}} 
\newcommand{\unitVal}{\codeText{unit}} 

\newcommand{\unitTyToty}{\unitTy{}{ }%
\codeText{$\rightarrow$ $\mathcal{T}$}}
\newcommand{\ty}{\codeText{$\mathcal{T}$}}

\newcommand{\tyOneToTyTwo}{\codeText{$\mathcal{T}_1$ $\rightarrow$ $\mathcal{T}_2$}}

\newcommand{\tyE}{\codeText{$E$}}
\newcommand{\tyTotyE}{\codeText{$\mathcal{T}$ $\rightarrow$ $E$}}
\newcommand{\tyToTyE}{\codeText{$\mathcal{T}$ $\rightarrow$ $E$}}
 
\newcommand{\tyValues}{\codeText{$\mathcal{T}$-}values}


\newcommand{\li}{\codeText{li}}

\newcommand{\tyNoun}{\codeText{Noun}}
\newcommand{\tyProposition}{\codeText{Proposition}}

\newcommand{\unicode}{\AcronymText{UNICODE}}
\newcommand{\AOP}{\AcronymText{AOP}} 

\newcommand{\chPkg}{\codeTextr{chPkg}}
\newcommand{\tCatChK}{\codeTextr{tCatChK}}

\newcommand{\sigmaTrans}{\codeText{sigmaTrans}} 

\newcommand{\chan}{\codeText{chan}}  

\newcommand{\ProcOne}{\ensuremath{\mathcal{P}_1}}


\newcommand{\carr}{\codeText{$\mathfrak{C}$}}
\newcommand{\tyEE}{\codeText{$\mathcal{E}$}}

\newcommand{\zero}{\codeText{0}} 
\newcommand{\szEqSizeFile}{\codeText{sz = size(file)}}

\newcommand{\sizeFn}{\codeText{size}} 

\newcommand{\openFileType}{\codeText{Open\_File}} 

\newcommand{\fileSym}{\codeText{file}} 
\newcommand{\szSym}{\codeText{sz}} 

\newcommand{\fileExistsFn}{\codeText{file\_exists}} 

\newcommand{\eli}{\item}


\newcommand{\yEqFx}{\codeText{y = f(x)}}

\newcommand{\bindFnSig}{\makebox{\codeText{($>>=$) :: m a $\Rightarrow$ %
(a $\rightarrow$ m b) $\rightarrow$ m b}}}

\renewcommand{\bindFn}{\codeText{$>>=$}}

\newcommand{\monadA}{\codeText{m a}}

\newcommand{\litOne}{\codeText{1}}
\newcommand{\litOnePtZ}{\codeText{1.0}}
\newcommand{\litOnePtO}{\codeText{1.1}}

\newcommand{\piOne}{\ensuremath{\pi_1}}
\newcommand{\piTwo}{\ensuremath{\pi_2}}
\newcommand{\piThree}{\ensuremath{\pi_3}}

\newcommand{\arrbl}{\ensuremath{\longrightarrow\hspace{-1.32em}\bullet}\hspace{.82em}}
\newcommand{\arrfo}{\ensuremath{\longrightarrow\hspace{-1.52em}\lhd}\hspace{.72em}}

\newcommand{\mGamma}{\ensuremath{\gamma}}

\newcommand{\pisbl}{\piOne{} \arrbl{} \piTwo{}}
\newcommand{\pisfo}{\piOne{} \arrfo{} $\ulcorner\piTwo{},\piThree{}\lrcorner$}



\newcommand{\OHO}{\codeText{101}}
\newcommand{\litOHO}{\codeText{101}}
\newcommand{\litOH}{\codeText{100}}
\newcommand{\litFive}{\codeText{5}}
\newcommand{\five}{\codeText{5}}

\newcommand{\nth}{{$n$}th}
%\newcommand{\nVal}{{$n$}}

\newcommand{\listval}{\codeText{list}}

\newcommand{\rlstsize}{\codeText{\ranged{0, list.size()-1}}}
\newcommand{\intrlstsize}{\codeText{int\ranged{0, list.size()-1}}}

\newcommand{\fIntG}{\codeText{int f(int g)}}
\newcommand{\fIntI}{\codeText{int f(int i)}}

\newcommand{\retFive}{\codeText{return 5}}

\newcommand{\QString}{\codeText{QString}}
\newcommand{\QObject}{\codeText{QObject}}
\newcommand{\ptrv}{\codeText{void*}}

\newcommand{\aVal}{\codeText{a}}

\newcommand{\switch}{\codeText{switch}}

\newcommand{\ztooh}{\codeText{0-100}}

\newcommand{\mOldLambda}{\ensuremath{\OldLambda}}



\renewenvironment{abstract}
 {\small
  \begin{center}
  \bfseries \abstractname\vspace{-.5em}\vspace{0pt}
  \end{center}
  \list{}{%
    \setlength{\leftmargin}{4mm}
    \setlength{\rightmargin}{\leftmargin}
  }
  \item\relax}
 {\endlist}


\newcommand{\pseudoIndent}{

\vspace{-2pt}\hspace{8pt}}
\newcommand{\pseudoInd}{\hspace{30pt}}

\newcommand{\caltypeT}{\ensuremath{\mathcal{T}}}
\newcommand{\calS}{\ensuremath{\mathcal{S}}}

\newcommand{\typeTp}{\ensuremath{\mathcal{T}'}}
\newcommand{\typeTpp}{\ensuremath{\mathcal{T}''}}

\setlist[itemize]{leftmargin=2mm}
\setlist[enumerate]{leftmargin=5mm}

\raggedbottom

\setcounter{secnumdepth}{3}

%\usepackage[capitalise,noabbrev,nameinlink]{cleveref}
%\Crefname{subsection}{Subsection}{Subsections}
%\crefname{subsection}{Subsection}{Subsections}
%\usepackage[none]{hyphenat}


